% =============================== 8. Q5 ===============================
\section{Question 5: What Competition Does to the Problem}\label{sec:q5}

\subsection{Three things competition can break}

The five questions were answered for a single operator holding the whole market.
Competition can affect availability, parking capacity and information, and the three are
not the same kind of claim: availability is a stochastic statement about stockouts,
parking capacity a deterministic one about slots, information a statement about what
any firm can estimate. We simulate the second at the Q4 plan's own resolution and treat
the first and third as arguments whose direction we can state, not whose magnitude we
can measure.

\subsection{A cell-resolution mismatch experiment}

The Q4 plan is fixed: 3,500 sites, 150 slots each, $R=450$ m, initial-inventory proxy
343,638. Two platforms share the market equally (share $s=0.5$), each drawing Poisson
arrivals and departures at half rate and holding half of every cell's inventory. Three
institutions are compared:
\begin{itemize}[leftmargin=1.2em,itemsep=0pt]
\item \textbf{independent} --- each platform owns a 150-slot fence at each open site and
  only its own bicycles park there;
\item \textbf{independent\_equal} --- the same, but 75 slots per platform, so the concrete
  per site equals the shared regime's 150 and the comparison isolates pooling, not total
  capacity;
\item \textbf{shared} --- one 150-slot fence per site, both platforms park in it, riders
  still take only their own bicycles, and overflow is rationed in proportion to that
  step's arrivals. Sharing pools slots, not fleets.
\end{itemize}
A spatial mismatch is then imposed at four intensities $m\in\{0,0.5,1.0,1.5\}$:
platform A's share of each origin cell's trips is $\sigma(a+2m\tilde x_i)$, where
$\sigma$ is the logistic function, $\tilde x_i$ is the cell's east--west coordinate scaled to $[-1,1]$ and $a$ is bisected
so that A still carries half of all trips. At $m=0$ the split is proportional and is the
model-consistent control; at $m=1.5$ A's riders are concentrated in the east while its
fleet remains at half of every cell's inventory. Twenty seeds, 15-minute steps,
about 1.48 million within-fence departures per run.

Every platform and site carries two states: bicycles parked \emph{inside} the fence,
$o_{jt}$, and bicycles parked just \emph{outside} it after an overflow, $u_{jt}$. A
departure draws from inside first and from outside second; only the remaining residual
is unmet. An arrival fills the inside room and the overflow enters $u_{jt}$, so
$o_{jt}+u_{jt}$ equals the platform's fleet at every step (the largest conservation
residual over all runs is $9\times10^{-11}$). Outside bicycles consume no slot and are
not reinserted into the
inside occupancy; the state exists so that no bicycle disappears from the accounting,
not to relax the parking constraint.

\begin{table}[htbp]
\centering
\footnotesize
\caption[Q5 spatial mismatch]{Two platforms at share 0.5 on the 3,500-site, 150-slot Q4
plan, by mismatch intensity $m$ and parking institution (20 seeds). ``Unmet'' is
unserved departures as a share of within-fence departures; ``full'' is the share of
arriving bicycles that overflowed a fence already at 150 slots; ``required slots'' is
the sum over sites of the 95th percentile over seeds of peak occupancy, an
unbounded-peak sizing convention, not existing berths.
The single-platform row is the same plan under one operator. Every row conserves
inside + outside = fleet.}
\label{tab:q5mismatch}
\scriptsize
\begin{tabular}{@{}llrrr@{}}
\toprule
$m$ & \textbf{Institution} & \textbf{Unmet rate} & \textbf{Full-site rate} &
\textbf{Required slots} \\
\midrule
--- & single platform & 17.15\% & 0.50\% & 479,675.68 \\
\midrule
0.0 & independent        & 17.80\% & 0.00002\% & 539,678.92 \\
0.0 & independent\_equal & 17.80\% & 0.93\%    & 539,678.92 \\
0.0 & shared             & 17.80\% & 0.44\%    & 503,073.56 \\
\midrule
0.5 & independent        & 19.65\% & 0.0052\%  & 535,780.02 \\
0.5 & independent\_equal & 19.65\% & 1.20\%    & 535,780.02 \\
0.5 & shared             & 19.65\% & 0.29\%    & 492,289.68 \\
\midrule
1.0 & independent        & 24.03\% & 0.0953\%  & 526,691.46 \\
1.0 & independent\_equal & 24.03\% & 2.04\%    & 526,691.46 \\
1.0 & shared             & 24.03\% & 0.29\%    & 470,548.04 \\
\midrule
1.5 & independent        & 29.50\% & 0.3123\%  & 515,894.53 \\
1.5 & independent\_equal & 29.50\% & 3.22\%    & 515,894.53 \\
1.5 & shared             & 29.50\% & 0.67\%    & 449,680.54 \\
\bottomrule
\end{tabular}
\end{table}

\subsection{What the simulation says}

\textbf{Splitting the market costs service even with no mismatch.} The single platform
leaves 17.15\% of within-fence departures unserved; at $m=0$, with both platforms
proportional, every two-platform institution leaves 17.80\%, about two thirds of a
percentage point worse. The mechanism is not spatial: each platform's Poisson
fluctuations no longer pool with the other's, so the same total fleet produces more
empty-fence steps. The size of the effect is specific to this plan and demand surface;
it is not a universal constant and does not carry across layouts. Twenty seeds bound
the noise: the seed-to-seed standard deviation of the unmet rate is
0.000728--0.000829 and of the full-site rate 0.000206--0.000255, so the 17.15\% against
17.80\% gap is about eight of those spreads. The shared and independent unmet rates, by
contrast, are equal in the raw output, not a statistical tie.

\textbf{The denominator is within-fence departures, not the city.} The 17.80\% is
263,307.06 unserved departures over 1,479,658.15 within-fence departures. It is
\emph{not} comparable with the 16.53\% departure-shortfall rate of the Q4 replay
(244,498.53 over 1,479,013.66 covered departures), because that figure is a single
operator on the 15-minute conserving replay while this one is two operators sharing
the same plan; whenever the two appear together they must be labelled separately.

\textbf{Mismatch is the larger effect.} Moving from $m=0$ to $m=1.5$ raises every
institution's unmet rate from 17.80\% to 29.50\%, nearly twelve percentage points. The
spatial concentration of one platform's riders away from its own bicycles is a stronger
force than the loss of noise pooling, and it grows monotonically with the intensity.

\textbf{Sharing fences does not reduce lost orders.} The shared regime's
unmet rate equals the independent regime's at every $m$ to the last digit (the
difference is 0.0 in the raw simulation output). The equality follows from the outside
state: availability depends on the fleet, not on the fence. Once inside is empty, a
rider draws from the outside pool; if both are empty the order is
lost regardless of who owns the concrete. Sharing pools parking slots, not bicycles.
What sharing does change is \emph{overflow}: the shared regime records 5,330 more
overflowing arrivals than independent at $m=0$ and 3,695 more at $m=1.5$, because a
shared fence has one wall while two independent fences have two, and the shared wall is
hit by both fleets. The extra overflow is absorbed by the outside state, which is why it
does not become an extra lost order.

\textbf{Sharing saves concrete.} The independent regime requires \textbf{1.0728} times
the shared capacity at $m=0$, 1.0883 at $m=0.5$, 1.1193 at $m=1.0$ and 1.1472 at
$m=1.5$: sharing the same fences would cut the slot requirement by \textbf{6.8\% to
12.8\%}, and the saving rises with mismatch. The shared regime needs 4.88\% more slots than a single platform at
$m=0$ and 6.25\% \emph{fewer} at $m=1.5$; the fall is a consequence of the service
failure, not a benefit, because fewer served departures mean fewer arrivals and lower
per-site peaks. These are required capacities --- a 95th-percentile per-site sizing
standard --- not existing berth counts, and they are specific to this 3,500-site,
150-slot, 450 m layout and this inventory; they do not transfer to another city or
another grid.

\textbf{The full-site rates separate the institutions cleanly.} The independent regime
almost never overflows (0.00002\% to 0.31\% of arriving bicycles): each platform's fence
sits half-empty while the other's riders walk past it. The independent\_equal regime
overflows up to 3.22\%, the shared regime sits between them at 0.29\% to 0.67\%, and the
single platform's 0.50\% exceeds independent's because its fences hold the whole fleet.
A low full-site rate is therefore evidence that capacity is in the wrong hands, not that
it is adequate.

\subsection{Cost and data duplication}

\textbf{Cost.} The rebalancing plan of Q2 is computed over the whole city; a platform
with half the trips still faces the same 290 arcs and the same hourly dispatch
window but fills each truck with half as many bicycles. The truck fleet's fixed
component --- 600 CNY
per truck per day, \textbf{42.96\%} of the daily cost at $K=60$ (36,000 of 83,807.74
CNY) --- is therefore incurred by every platform separately, and the same holds for
the tricycle layer's fixed cost, 300 CNY per tricycle-day, which is charged on each
platform's own fleet.

\textbf{Data.} Every model here is
estimated from the full origin--destination matrix, and a platform holding half of
the market observes half of the rows, so the attraction $A_j$ of
\eqnref{eq:gravity} and the cell deficit $\Delta_c$ of \eqnref{eq:q3lp} are estimated
with a standard error that grows as $\sqrt{2}$. More decisively, \emph{no platform can
compute the plans of Q2 and Q4 at all}: the Q2 truck plan is defined on a city-wide
zone network, the Q4 layout on a city-wide demand surface, and the Q3 reward on a
city-wide deficit list. Each is a public-good computation that only the city can run.

\subsection{Multi-homing, and why it does not rescue the situation}\label{sec:multihome}

If riders use several applications, availability may improve: a rider who finds no
bicycle on one platform may find one on another, so effective density is closer to the
sum of the platforms' densities. But multi-homing shifts competition to a different
margin: platforms compete for the same trips, so each must offer near-complete coverage
to be usable, and the fleet that results is the sum of $n$ fleets each sized for much of
the market rather than $1/n$ of it. Multi-homing therefore trades a possible service
benefit for a capital cost, and the capital cost is the larger of the two in a
saturated market.

Neither regime escapes the underlying externality: a deployed bicycle
occupies public space and adds to the stock all platforms draw on, but the
platform bears only its own deployment cost. This is the standard common-pool
argument \cite{katz1985,rochet2003}, and in a market where the marginal bicycle adds
almost no service (\S\ref{sec:q2}) it predicts what was observed in 2017:
fleets growing faster than ridership, and consolidation once the capital could no
longer be financed.

\subsection{What a regulator should do}

\textbf{Share the data, not the bicycles}: a city-wide OD matrix is
the binding public good, and requiring platforms to deposit trip records with a
neutral trustee costs the firms nothing competitively. \textbf{Share the fences, and
pool more than concrete}: the shared regime saves 6.8\%--12.8\% of the slot
requirement, but shared fences alone do not reduce lost orders (\S\ref{sec:q5}),
because the fleets remain separate. \textbf{Coordinate rebalancing}: a neutral
rebalancing operator, or a requirement that platforms' trucks serve shared depots,
recovers a fixed cost currently paid twice. \textbf{Tie fleet caps to measured
service}: a cap expressed as bicycles per thousand trips, rather than as a fixed
number, stays anchored to the measured service requirement and is auditable from the
shared data of the first measure.

\subsection{What the external data say about the market}

The OpenStreetMap check of \S\ref{sec:q4} is also evidence about competition. The
attractors the model identifies --- the busiest destination cells --- are public
places: subway entrances, commercial streets, university gates, not platform-specific.
The city is the natural owner of the model: data sharing is an argument about who can see the city
whole.
